% ============================================================
%  Section 3 – Choix méthodologiques
% ============================================================
%
%  Contenu attendu (voir sujet PDF) :
%
%  3.1 Outils d'inférence retenus
%      - Justifier le choix entre IC de proportion (Déf. 2.11) et
%        IC de la moyenne (Déf. 2.14, loi de Student)
%      - Justifier le choix entre test de conformité proportion / moyenne
%        et test du χ² de conformité / d'indépendance
%
%  3.2 Justification des conditions de validité
%      - Pour chaque test : rappeler les conditions (n≥30, np≥5, n≥50…)
%      - Vérifier explicitement sur les données (donner les valeurs)
%      - Si conditions non réunies : le signaler et présenter quand même
%        le test, conformément à la consigne du professeur
%
%  3.3 Choix loi normale vs loi de Student
%      - Rappeler : σ inconnu → Student T(n−1) ; σ connu → N(0,1)
%      - Dans notre cas : variance estimée Se² = n/(n−1) × σ²obs (cours § 2.3.1)
%
% ============================================================

\section{Choix méthodologiques}

\subsection{Outils d'inférence retenus}

% TODO : rédiger ici

\subsection{Vérification des conditions de validité}

% TODO : tableau : test | condition | valeur observée | validée ?

\subsection{Justification loi normale vs loi de Student}

% TODO : rédiger ici
